\documentclass{article}
\usepackage[utf8]{inputenc}
\usepackage{amsmath}
\usepackage{amssymb}
\usepackage{amsthm}
\usepackage{geometry}
\usepackage{booktabs}
\usepackage{hyperref}

\geometry{a4paper, margin=1in}

\title{Summary of Previous Tested A* Update Strategies and Cost Functions}
\date{}
\begin{document}

\maketitle


\section{A* Search Cost Components}

The A* algorithm evaluates nodes $n$ in the search space using a total estimated cost function $f(n)$, which is typically a combination of the cost-to-come $g(n)$ and the heuristic cost-to-go $h(n)$.

\subsection{Heuristic Cost (Cost-to-Go), $h(n)$}
The heuristic $h(n)$ represents the estimated cost from the current state (node $n$) to the goal state. In this context, $h(n)$ is defined as the model's loss on the training data $(X, Y)$:
$$h(n) = \text{Loss}(\text{MLP}_n(X), Y)$$
where $\text{MLP}_n$ is the Quantized MLP model associated with node $n$, $\text{Loss}$ is the chosen loss function and $T_L$ is the target loss.

\subsection{Cost-to-Come, $g(n)$}
The cost-to-come $g(n)$ is the accumulated actual cost from the initial state $n_0$ to the current state $n$. For a transition from a parent node $n_{parent}$ to a neighbor node $n$, the cost is updated by a step cost $\Delta g$:
$$g(n) = g(n_{parent}) + \Delta g(n_{parent} \to n)$$
The initial cost is $g(n_0) = g_{\text{ini}}$ and the step cost $\Delta g$ is dependent on the chosen update strategy.

\subsection{Total Estimated Cost, $f(n)$}
The total estimated cost $f(n)$ is calculated based on $g(n)$ and $h(n)$, and may incorporate a scaling factor to dynamically adjust the influence of $g$ versus $h$.

\subsubsection{Standard A* (Unscaled $f$)}
If $\text{update\_strategy} = 0$ the standard A* function is used:
$$f(n) = g(n) + h(n)$$

\subsubsection{Scaled $f$ (for Strategies 1, 2 and 3)}
If a scaling factor is applied, the cost-to-come $g(n)$ is scaled by a term that depends on the current loss $h(n)$ relative to the target loss $T_L$. This essentially prioritizes nodes with a loss close to the target.
$$f(n) = \max\left(0, 1 - \frac{T_L}{h(n)}\right) g(n) + h(n)$$

\section{Update Strategies for Step Cost $\Delta g$}

The step cost $\Delta g(n_{parent} \to n)$ for moving from a parent node $n_{parent}$ to a neighbor node $n$ depends on the parameter $\text{update\_strategy}$.

Let $h_{parent} = h(n_{parent})$ and $h = h(n)$. Let $h_0$ be the loss of the initial node.
\begin{itemize}
    \item $C$: Total number of possible quantized weight configurations.
    \item $M$: Maximum number of search iterations.
    \item $\alpha$: Mixing factor ($0 \le \alpha \le 1$).
    \item $g_{\text{step}}$: A constant step value.
    \item $k$: Current iteration index.
\end{itemize}

\subsection{Strategy 0: Constant Step Size}
The step cost is a fixed parameter:
$$\Delta g = g_{\text{step}}$$

\subsection{Strategy 1: Loss-Based Scaling (Log-Scale Configuration)}
The step cost is scaled based on the relative change in loss from the parent node and the initial node.
$$\Delta g = \frac{1}{\ln(C)} \left( \alpha \frac{h}{h_{parent}} + (1-\alpha) \frac{h}{h_0} \right)$$

\subsection{Strategy 2: Normalized Constant Step Size}
The step cost is a small, constant fraction of the maximum allowed iterations:
$$\Delta g = \frac{1}{M}$$

\subsection{Strategy 3: Loss-Based Scaling with Iteration-Dependent Exponent}
This strategy uses a complex scaling factor that is further modulated by the iteration number $k$ using a base-10 logarithm.
$$\Delta g = \frac{1}{M \cdot \log_{10}(C)} \left( \alpha \frac{h}{h_{parent}} + (1-\alpha) \frac{h}{h_0} \right)^{\max(1, \log_{10}(k))}$$

\section{State Space}

The search space is discrete, defined by the quantization of the model parameters.

\subsection{Quantization Step}
The quantization step $\delta$ is determined by the quantization factor $Q = \text{quantization\_factor}$:
$$\delta = \frac{1}{Q}$$
A weight $w$ is quantized to $w_q$ by rounding to the nearest multiple of $\delta$:
$$w_q = \frac{\text{round}(w \cdot Q)}{Q}$$

\subsection{Total Possible Configurations}
The total number of possible quantized configurations, $C$, is used to normalize the step cost in Strategies 1 and 3. Let $P$ be the total number of network parameters and $r_{\max} = \text{parameter\_range}[1]$ be the maximum absolute parameter value (assuming a symmetric range). The number of possible values for a single parameter is $2 \cdot r_{\max} \cdot Q + 1$.
$$C = (2 \cdot r_{\max} \cdot Q + 1)^P$$

\end{document}